
\documentclass{article}
\usepackage{colortbl}
\usepackage{makecell}
\usepackage{multirow}
\usepackage{supertabular}

\begin{document}

\newcounter{utterance}

\centering \large Interaction Transcript for game `hot\_air\_balloon', experiment `air\_balloon\_survival\_de\_negotiation\_easy', episode 2 with deepseek{-}chat{-}v3.1{-}t1.0.
\vspace{24pt}

{ \footnotesize  \setcounter{utterance}{1}
\setlength{\tabcolsep}{0pt}
\begin{supertabular}{c@{$\;$}|p{.15\linewidth}@{}p{.15\linewidth}p{.15\linewidth}p{.15\linewidth}p{.15\linewidth}p{.15\linewidth}}
    \# & $\;$A & \multicolumn{4}{c}{Game Master} & $\;\:$B\\
    \hline

    \theutterance \stepcounter{utterance}  
    & & \multicolumn{4}{p{0.6\linewidth}}{
        \cellcolor[rgb]{0.9,0.9,0.9}{
            \makecell[{{p{\linewidth}}}]{
                \texttt{\tiny{[P1$\langle$GM]}}
                \texttt{Du nimmst an einem kollaborativen Verhandlungspiel Teil.} \\
\\ 
\texttt{Zusammen mit einem anderen Teilnehmer musst du dich auf eine Reihe von Gegenständen entscheiden, die behalten werden. Jeder von euch hat eine persönliche Verteilung über die Wichtigkeit der einzelnen Gegenstände. Jeder von euch hat eine eigene Meinung darüber, wie wichtig jeder einzelne Gegenstand ist (Gegenstandswichtigkeit). Du kennst die Wichtigkeitsverteilung des anderen Spielers nicht. Zusätzlich siehst du, wie viel Aufwand jeder Gegenstand verursacht.  } \\
\texttt{Ihr dürft euch nur auf eine Reihe von Gegenständen einigen, wenn der Gesamtaufwand der ausgewählten Gegenstände den Maximalaufwand nicht überschreitet:} \\
\\ 
\texttt{Maximalaufwand = 3288} \\
\\ 
\texttt{Hier sind die einzelnen Aufwände der Gegenstände:} \\
\\ 
\texttt{Aufwand der Gegenstände = \{"A07": 329, "A36": 41, "B73": 417, "C45": 75, "B93": 390, "B98": 887, "A87": 808, "A63": 151, "B05": 849, "C13": 129, "B16": 350, "B33": 118, "C34": 630, "C89": 602, "C11": 801\}} \\
\\ 
\texttt{Hier ist deine persönliche Verteilung der Wichtigkeit der einzelnen Gegenstände:} \\
\\ 
\texttt{Werte der Gegenstandswichtigkeit = \{"A07": 138, "A36": 583, "B73": 868, "C45": 822, "B93": 783, "B98": 65, "A87": 262, "A63": 121, "B05": 508, "C13": 780, "B16": 461, "B33": 484, "C34": 668, "C89": 389, "C11": 808\}} \\
\\ 
\texttt{Ziel:} \\
\\ 
\texttt{Dein Ziel ist es, eine Reihe von Gegenständen auszuhandeln, die dir möglichst viel bringt (d. h. Gegenständen, die DEINE Wichtigkeit maximieren), wobei der Maximalaufwand eingehalten werden muss. Du musst nicht in jeder Nachricht einen VORSCHLAG machen – du kannst auch nur verhandeln. Alle Taktiken sind erlaubt!} \\
\\ 
\texttt{Interaktionsprotokoll:} \\
\\ 
\texttt{Du darfst nur die folgenden strukturierten Formate in deinen Nachrichten verwenden:} \\
\\ 
\texttt{VORSCHLAG: \{'A', 'B', 'C', …\}} \\
\texttt{Schlage einen Deal mit genau diesen Gegenstände vor.} \\
\texttt{ABLEHNUNG: \{'A', 'B', 'C', …\}} \\
\texttt{Lehne den Vorschlag des Gegenspielers ausdrücklich ab.} \\
\texttt{ARGUMENT: \{'...'\}} \\
\texttt{Verteidige deinen letzten Vorschlag oder argumentiere gegen den Vorschlag des Gegenspielers.} \\
\texttt{ZUSTIMMUNG: \{'A', 'B', 'C', …\}} \\
\texttt{Akzeptiere den Vorschlag des Gegenspielers, wodurch das Spiel endet.} \\
\texttt{STRATEGISCHE ÜBERLEGUNGEN: \{'...'\}} \\
\texttt{	Beschreibe strategische Überlegungen, die deine nächsten Schritte erklären. Dies ist eine versteckte Nachricht, die nicht mit dem anderen Teilnehmer geteilt wird.} \\
\\ 
\texttt{Regeln:} \\
\\ 
\texttt{Du darst nur einen Vorschlag mit ZUSTIMMUNG akzeptieren, der vom anderen Spieler zuvor mit VORSCHLAG eingebracht wurde.} \\
\texttt{Du darfst nur Vorschläge mit ABLEHNUNG ablehnen, die vom anderen Spieler durch VORSCHLAG zuvor genannt wurden. } \\
\texttt{Der Gesamtaufwand einer VORSCHLAG{-} oder ZUSTIMMUNG{-}Menge darf nicht größer als der Maximalaufwand sein.  } \\
\texttt{Offenbare deine versteckte Wichtigkeitsverteilung nicht.} \\
\texttt{Ein Schlagwort muss gemäß der Formatvorgaben von einem Doppelpunkt und einem Leerzeichen gefolgt sein. Das Argument muss eine Python{-}Menge mit 0 oder mehr Strings sein.  } \\
\texttt{Die Form muss also SCHLAGWORT: \{...\} sein.} \\
\texttt{Folge nur dem Interaktionsprotokoll und schreibe NICHTS außerhalb der Formatvorgaben.} \\
\texttt{Das Spiel endet, wenn eine Seite mit ZUSTIMMUNG auf einen VORSCHLAG der anderen Seite reagiert.  } \\
\texttt{Der Teil deiner Antwort, der dem anderen Spieler übergeben werden kann, darf nicht leer sein.  } \\
\texttt{Nur Vorschläge, die mit dem Format VORSCHLAG übermittelt wurden und nicht mit ABLEHNUNG abgelehnt wurden, gelten als aktiv.  } \\
\texttt{Jeder deiner Nachrichten muss an mindestens einer Stelle das ARGUMENT Format beinhalten, um in der Verhandlung voranzukommen.} \\
\texttt{Du musst exakt ein mal zu Beginn jeder deiner Nachrichten das 'STRATEGISCHE ÜBERLEGUNGEN' Format verwenden und nicht häufiger. Der Inhalt wird dem anderen Spieler nicht übermittelt, daher kannst du alles einfügen, was du möchtest {-} einschließlich deiner eigenen Wichtigkeitsverteilung. Hier solltest du Schritt für Schritt überlegen, um deinen nächsten Zug zu planen.} \\
            }
        }
    }
    & \\ \\

    \theutterance \stepcounter{utterance}  
    & \multicolumn{4}{p{0.6\linewidth}}{
        \cellcolor[rgb]{0.9,0.9,0.9}{
            \makecell[{{p{\linewidth}}}]{
                \texttt{\tiny{[P1$\rangle$GM]}}
                \texttt{STRATEGISCHE ÜBERLEGUNGEN: \{'Ich muss meine hohen Wertgegenstände mit niedrigem Aufwand priorisieren: A36 (583/41), C45 (822/75), B73 (868/417), C13 (780/129). Meine Top{-}Performer sind A36 (Wert/Aufwand ≈ 14.22) und C45 (≈10.96). Maximalaufwand ist 3288. Ich sollte zunächst meine starke Position vorschlagen und dann verhandeln.'\}} \\
\texttt{ARGUMENT: \{'Ich schlage vor, dass wir uns auf Gegenstände konzentrieren, die einen hohen Wert bei relativ geringem Aufwand bieten, um die Gesamtressourcen optimal zu nutzen.'\}} \\
\texttt{VORSCHLAG: \{'A36', 'C45', 'B73', 'C13', 'B93', 'B16', 'B33', 'C34', 'C89', 'C11'\}} \\
            }
        }
    }
    & & \\ \\

    \theutterance \stepcounter{utterance}  
    & & \multicolumn{4}{p{0.6\linewidth}}{
        \cellcolor[rgb]{0.9,0.9,0.9}{
            \makecell[{{p{\linewidth}}}]{
                \texttt{\tiny{[GM$\rangle$P2]}}
                \texttt{Du nimmst an einem kollaborativen Verhandlungspiel Teil.} \\
\\ 
\texttt{Zusammen mit einem anderen Teilnehmer musst du dich auf eine Reihe von Gegenständen entscheiden, die behalten werden. Jeder von euch hat eine persönliche Verteilung über die Wichtigkeit der einzelnen Gegenstände. Jeder von euch hat eine eigene Meinung darüber, wie wichtig jeder einzelne Gegenstand ist (Gegenstandswichtigkeit). Du kennst die Wichtigkeitsverteilung des anderen Spielers nicht. Zusätzlich siehst du, wie viel Aufwand jeder Gegenstand verursacht.  } \\
\texttt{Ihr dürft euch nur auf eine Reihe von Gegenständen einigen, wenn der Gesamtaufwand der ausgewählten Gegenstände den Maximalaufwand nicht überschreitet:} \\
\\ 
\texttt{Maximalaufwand = 3288} \\
\\ 
\texttt{Hier sind die einzelnen Aufwände der Gegenstände:} \\
\\ 
\texttt{Aufwand der Gegenstände = \{"A07": 329, "A36": 41, "B73": 417, "C45": 75, "B93": 390, "B98": 887, "A87": 808, "A63": 151, "B05": 849, "C13": 129, "B16": 350, "B33": 118, "C34": 630, "C89": 602, "C11": 801\}} \\
\\ 
\texttt{Hier ist deine persönliche Verteilung der Wichtigkeit der einzelnen Gegenstände:} \\
\\ 
\texttt{Werte der Gegenstandswichtigkeit = \{"A07": 138, "A36": 583, "B73": 868, "C45": 822, "B93": 783, "B98": 65, "A87": 262, "A63": 121, "B05": 508, "C13": 780, "B16": 461, "B33": 484, "C34": 668, "C89": 389, "C11": 808\}} \\
\\ 
\texttt{Ziel:} \\
\\ 
\texttt{Dein Ziel ist es, eine Reihe von Gegenständen auszuhandeln, die dir möglichst viel bringt (d. h. Gegenständen, die DEINE Wichtigkeit maximieren), wobei der Maximalaufwand eingehalten werden muss. Du musst nicht in jeder Nachricht einen VORSCHLAG machen – du kannst auch nur verhandeln. Alle Taktiken sind erlaubt!} \\
\\ 
\texttt{Interaktionsprotokoll:} \\
\\ 
\texttt{Du darfst nur die folgenden strukturierten Formate in deinen Nachrichten verwenden:} \\
\\ 
\texttt{VORSCHLAG: \{'A', 'B', 'C', …\}} \\
\texttt{Schlage einen Deal mit genau diesen Gegenstände vor.} \\
\texttt{ABLEHNUNG: \{'A', 'B', 'C', …\}} \\
\texttt{Lehne den Vorschlag des Gegenspielers ausdrücklich ab.} \\
\texttt{ARGUMENT: \{'...'\}} \\
\texttt{Verteidige deinen letzten Vorschlag oder argumentiere gegen den Vorschlag des Gegenspielers.} \\
\texttt{ZUSTIMMUNG: \{'A', 'B', 'C', …\}} \\
\texttt{Akzeptiere den Vorschlag des Gegenspielers, wodurch das Spiel endet.} \\
\texttt{STRATEGISCHE ÜBERLEGUNGEN: \{'...'\}} \\
\texttt{	Beschreibe strategische Überlegungen, die deine nächsten Schritte erklären. Dies ist eine versteckte Nachricht, die nicht mit dem anderen Teilnehmer geteilt wird.} \\
\\ 
\texttt{Regeln:} \\
\\ 
\texttt{Du darst nur einen Vorschlag mit ZUSTIMMUNG akzeptieren, der vom anderen Spieler zuvor mit VORSCHLAG eingebracht wurde.} \\
\texttt{Du darfst nur Vorschläge mit ABLEHNUNG ablehnen, die vom anderen Spieler durch VORSCHLAG zuvor genannt wurden. } \\
\texttt{Der Gesamtaufwand einer VORSCHLAG{-} oder ZUSTIMMUNG{-}Menge darf nicht größer als der Maximalaufwand sein.  } \\
\texttt{Offenbare deine versteckte Wichtigkeitsverteilung nicht.} \\
\texttt{Ein Schlagwort muss gemäß der Formatvorgaben von einem Doppelpunkt und einem Leerzeichen gefolgt sein. Das Argument muss eine Python{-}Menge mit 0 oder mehr Strings sein.  } \\
\texttt{Die Form muss also SCHLAGWORT: \{...\} sein.} \\
\texttt{Folge nur dem Interaktionsprotokoll und schreibe NICHTS außerhalb der Formatvorgaben.} \\
\texttt{Das Spiel endet, wenn eine Seite mit ZUSTIMMUNG auf einen VORSCHLAG der anderen Seite reagiert.  } \\
\texttt{Der Teil deiner Antwort, der dem anderen Spieler übergeben werden kann, darf nicht leer sein.  } \\
\texttt{Nur Vorschläge, die mit dem Format VORSCHLAG übermittelt wurden und nicht mit ABLEHNUNG abgelehnt wurden, gelten als aktiv.  } \\
\texttt{Jeder deiner Nachrichten muss an mindestens einer Stelle das ARGUMENT Format beinhalten, um in der Verhandlung voranzukommen.} \\
\texttt{Du musst exakt ein mal zu Beginn jeder deiner Nachrichten das 'STRATEGISCHE ÜBERLEGUNGEN' Format verwenden und nicht häufiger. Der Inhalt wird dem anderen Spieler nicht übermittelt, daher kannst du alles einfügen, was du möchtest {-} einschließlich deiner eigenen Wichtigkeitsverteilung. Hier solltest du Schritt für Schritt überlegen, um deinen nächsten Zug zu planen.} \\
\\ 
\texttt{Du erhältst nun die erste Nachricht des anderen Spielers.} \\
\\ 
\texttt{ARGUMENT: \{'Ich schlage vor, dass wir uns auf Gegenstände konzentrieren, die einen hohen Wert bei relativ geringem Aufwand bieten, um die Gesamtressourcen optimal zu nutzen.'\}} \\
\texttt{VORSCHLAG: \{'A36', 'C45', 'B73', 'C13', 'B93', 'B16', 'B33', 'C34', 'C89', 'C11'\}} \\
            }
        }
    }
    & \\ \\

    \theutterance \stepcounter{utterance}  
    & & & \multicolumn{4}{p{0.6\linewidth}}{
        \cellcolor[rgb]{0.9,0.9,0.9}{
            \makecell[{{p{\linewidth}}}]{
                \texttt{\tiny{[GM$\langle$P2]}}
                \texttt{STRATEGISCHE ÜBERLEGUNGEN: \{'Der Gegner hat einen Vorschlag gemacht, der 10 Gegenstände enthält. Ich muss den Gesamtaufwand und meinen eigenen Wert berechnen. Mein Ziel ist es, meinen Gesamtwert zu maximieren, während der Maximalaufwand von 3288 eingehalten wird. Zuerst berechne ich den Aufwand und Wert des vorgeschlagenen Deals vom Gegner.'\}} \\
\texttt{ARGUMENT: \{'Dein Vorschlag hat einen Gesamtaufwand von: A36(41) + C45(75) + B73(417) + C13(129) + B93(390) + B16(350) + B33(118) + C34(630) + C89(602) + C11(801) = 41+75+417+129+390+350+118+630+602+801 = 3553. Dieser übersteigt den Maximalaufwand von 3288, ist also ungültig und kann nicht akzeptiert werden. Wir müssen einen Deal finden, der innerhalb des Budgets liegt.'\}} \\
            }
        }
    }
     \\ \\

    \theutterance \stepcounter{utterance}  
    & & \multicolumn{4}{p{0.6\linewidth}}{
        \cellcolor[rgb]{0.9,0.9,0.9}{
            \makecell[{{p{\linewidth}}}]{
                \texttt{\tiny{[P1$\langle$GM]}}
                \texttt{ARGUMENT: \{'Dein Vorschlag hat einen Gesamtaufwand von: A36(41) + C45(75) + B73(417) + C13(129) + B93(390) + B16(350) + B33(118) + C34(630) + C89(602) + C11(801) = 41+75+417+129+390+350+118+630+602+801 = 3553. Dieser übersteigt den Maximalaufwand von 3288, ist also ungültig und kann nicht akzeptiert werden. Wir müssen einen Deal finden, der innerhalb des Budgets liegt.'\}} \\
            }
        }
    }
    & \\ \\

    \theutterance \stepcounter{utterance}  
    & \multicolumn{4}{p{0.6\linewidth}}{
        \cellcolor[rgb]{0.9,0.9,0.9}{
            \makecell[{{p{\linewidth}}}]{
                \texttt{\tiny{[P1$\rangle$GM]}}
                \texttt{STRATEGISCHE ÜBERLEGUNGEN: \{'Mein Vorschlag war zu großzügig und hat das Budget überschritten. Ich muss einige Gegenstände entfernen, um unter 3288 zu bleiben. Ich sollte meine Top{-}Performer behalten: A36 (583/41), C45 (822/75), B73 (868/417), C13 (780/129). Diese haben zusammen einen Aufwand von 41+75+417+129 = 662 und einen Wert von 583+822+868+780 = 3053. Ich kann noch weitere wertvolle Gegenstände hinzufügen, die in das Budget passen.'\}} \\
\texttt{ARGUMENT: \{'Du hast recht, mein erster Vorschlag war über dem Budget. Lassen wir uns auf eine Auswahl konzentrieren, die unseren gemeinsamen Maximalaufwand nicht überschreitet, aber dennoch für beide Seiten wertvoll ist.'\}} \\
\texttt{VORSCHLAG: \{'A36', 'C45', 'B73', 'C13', 'B93', 'B16', 'B33', 'C34'\}} \\
            }
        }
    }
    & & \\ \\

    \theutterance \stepcounter{utterance}  
    & & \multicolumn{4}{p{0.6\linewidth}}{
        \cellcolor[rgb]{0.9,0.9,0.9}{
            \makecell[{{p{\linewidth}}}]{
                \texttt{\tiny{[GM$\rangle$P2]}}
                \texttt{ARGUMENT: \{'Du hast recht, mein erster Vorschlag war über dem Budget. Lassen wir uns auf eine Auswahl konzentrieren, die unseren gemeinsamen Maximalaufwand nicht überschreitet, aber dennoch für beide Seiten wertvoll ist.'\}} \\
\texttt{VORSCHLAG: \{'A36', 'C45', 'B73', 'C13', 'B93', 'B16', 'B33', 'C34'\}} \\
            }
        }
    }
    & \\ \\

    \theutterance \stepcounter{utterance}  
    & & & \multicolumn{4}{p{0.6\linewidth}}{
        \cellcolor[rgb]{0.9,0.9,0.9}{
            \makecell[{{p{\linewidth}}}]{
                \texttt{\tiny{[GM$\langle$P2]}}
                \texttt{STRATEGISCHE ÜBERLEGUNGEN: \{'Der Gegner schlägt jetzt 8 Gegenstände vor: A36, C45, B73, C13, B93, B16, B33, C34. Ich berechne den Aufwand: 41+75+417+129+390+350+118+630 = 2150. Das ist weit unter dem Budget von 3288. Mein Wert dafür: 583+822+868+780+783+461+484+668 = 5449. Das ist gut, aber ich könnte vielleicht mehr Wert innerhalb des Budgets unterbringen. Ich sollte Gegenstände mit hohem Wert{-}Aufwand{-}Verhältnis hinzufügen, insbesondere B05 (508 Wert, 849 Aufwand) und C11 (808 Wert, 801 Aufwand) haben hohen absoluten Wert. A07 (138 Wert, 329 Aufwand) und A63 (121 Wert, 151 Aufwand) sind weniger effizient. Ich schlage einen counter deal vor, der mehr Wert für mich bringt.'\}} \\
\texttt{ARGUMENT: \{'Danke für den angepassten Vorschlag. Während er innerhalb des Budgets liegt, glaube ich, dass wir die Gesamtwerte weiter optimieren können, indem wir einige zusätzliche hochwertige Gegenstände hinzufügen, ohne das Limit zu überschreiten.'\}} \\
\texttt{VORSCHLAG: \{'A36', 'C45', 'B73', 'C13', 'B93', 'B16', 'B33', 'C34', 'B05', 'C11'\}} \\
            }
        }
    }
     \\ \\

    \theutterance \stepcounter{utterance}  
    & & \multicolumn{4}{p{0.6\linewidth}}{
        \cellcolor[rgb]{0.9,0.9,0.9}{
            \makecell[{{p{\linewidth}}}]{
                \texttt{\tiny{[P1$\langle$GM]}}
                \texttt{ARGUMENT: \{'Danke für den angepassten Vorschlag. Während er innerhalb des Budgets liegt, glaube ich, dass wir die Gesamtwerte weiter optimieren können, indem wir einige zusätzliche hochwertige Gegenstände hinzufügen, ohne das Limit zu überschreiten.'\}} \\
\texttt{VORSCHLAG: \{'A36', 'C45', 'B73', 'C13', 'B93', 'B16', 'B33', 'C34', 'B05', 'C11'\}} \\
            }
        }
    }
    & \\ \\

    \theutterance \stepcounter{utterance}  
    & \multicolumn{4}{p{0.6\linewidth}}{
        \cellcolor[rgb]{0.9,0.9,0.9}{
            \makecell[{{p{\linewidth}}}]{
                \texttt{\tiny{[P1$\rangle$GM]}}
                \texttt{STRATEGISCHE ÜBERLEGUNGEN: \{'Der Gegner schlägt vor, B05 (849 Aufwand) und C11 (801 Aufwand) hinzuzufügen. Mein aktueller Vorschlag hat einen Aufwand von: A36(41) + C45(75) + B73(417) + C13(129) + B93(390) + B16(350) + B33(118) + C34(630) = 2150. Mit B05 und C11 wäre der Aufwand 2150 + 849 + 801 = 3800, was weit über 3288 liegt. Das ist nicht machbar. Ich sollte einen Gegenvorschlag machen, der nah an 3288 liegt, aber meine Top{-}Performer behält. Vielleicht B05 statt C34? C34 hat für mich einen Wert von 668 bei 630 Aufwand, B05 hat 508 bei 849 Aufwand {-} das wäre ein schlechter Tausch für mich. Lieber C11 (808 Wert/801 Aufwand) statt C34 (668/630) anbieten? Das gibt mir mehr Wert, aber auch mehr Aufwand. Ich sollte zuerst den Aufwand des Gegenvorschlags prüfen.'\}} \\
\\ 
\texttt{ARGUMENT: \{'Dein Vorschlag überschreitet das Budget deutlich. Lassen Sie uns stattdessen meinen Vorschlag als Basis nehmen und vielleicht einen zusätzlichen Gegenstand hinzufügen, der für beide wertvoll ist, aber innerhalb des Limits bleibt.'\}} \\
\texttt{VORSCHLAG: \{'A36', 'C45', 'B73', 'C13', 'B93', 'B16', 'B33', 'C34', 'C89'\}} \\
            }
        }
    }
    & & \\ \\

    \theutterance \stepcounter{utterance}  
    & & \multicolumn{4}{p{0.6\linewidth}}{
        \cellcolor[rgb]{0.9,0.9,0.9}{
            \makecell[{{p{\linewidth}}}]{
                \texttt{\tiny{[GM$\rangle$P2]}}
                \texttt{ARGUMENT: \{'Dein Vorschlag überschreitet das Budget deutlich. Lassen Sie uns stattdessen meinen Vorschlag als Basis nehmen und vielleicht einen zusätzlichen Gegenstand hinzufügen, der für beide wertvoll ist, aber innerhalb des Limits bleibt.'\}} \\
\texttt{VORSCHLAG: \{'A36', 'C45', 'B73', 'C13', 'B93', 'B16', 'B33', 'C34', 'C89'\}} \\
            }
        }
    }
    & \\ \\

    \theutterance \stepcounter{utterance}  
    & & & \multicolumn{4}{p{0.6\linewidth}}{
        \cellcolor[rgb]{0.9,0.9,0.9}{
            \makecell[{{p{\linewidth}}}]{
                \texttt{\tiny{[GM$\langle$P2]}}
                \texttt{STRATEGISCHE ÜBERLEGUNGEN: \{'Der Gegner schlägt seinen vorherigen Deal plus C89 vor. Aufwand: 2150 (von vorher) + 602 (C89) = 2752. Das ist unter 3288. Mein Wert: 5449 (von vorher) + 389 (C89) = 5838. Das ist gut, aber ich könnte immer noch B05 (Wert 508, Aufwand 849) hinzufügen, was 2752 + 849 = 3601 ergeben würde, was über dem Budget liegt. Vielleicht kann ich einen weniger wertvollen Gegenstand entfernen, um Platz für B05 zu schaffen. A36 hat einen hohen Wert für mich (583) aber niedrigen Aufwand (41), also will ich das behalten. Vielleicht entferne ich C89 (Wert 389, Aufwand 602) und füge B05 hinzu: Aufwand 2752 {-} 602 + 849 = 2999, Wert 5838 {-} 389 + 508 = 5957. Das ist besser für mich. Ich schlage einen counter deal vor.'\}} \\
\texttt{ARGUMENT: \{'Dein Vorschlag mit C89 ist solide, aber ich denke, wir können einen noch besseren Gesamtwert erzielen, indem wir C89 gegen einen anderen wertvollen Gegenstand austauschen, der besser in das Budget passt.'\}} \\
\texttt{VORSCHLAG: \{'A36', 'C45', 'B73', 'C13', 'B93', 'B16', 'B33', 'C34', 'B05'\}} \\
            }
        }
    }
     \\ \\

    \theutterance \stepcounter{utterance}  
    & & \multicolumn{4}{p{0.6\linewidth}}{
        \cellcolor[rgb]{0.9,0.9,0.9}{
            \makecell[{{p{\linewidth}}}]{
                \texttt{\tiny{[P1$\langle$GM]}}
                \texttt{ARGUMENT: \{'Dein Vorschlag mit C89 ist solide, aber ich denke, wir können einen noch besseren Gesamtwert erzielen, indem wir C89 gegen einen anderen wertvollen Gegenstand austauschen, der besser in das Budget passt.'\}} \\
\texttt{VORSCHLAG: \{'A36', 'C45', 'B73', 'C13', 'B93', 'B16', 'B33', 'C34', 'B05'\}} \\
            }
        }
    }
    & \\ \\

    \theutterance \stepcounter{utterance}  
    & \multicolumn{4}{p{0.6\linewidth}}{
        \cellcolor[rgb]{0.9,0.9,0.9}{
            \makecell[{{p{\linewidth}}}]{
                \texttt{\tiny{[P1$\rangle$GM]}}
                \texttt{STRATEGISCHE ÜBERLEGUNGEN: \{'Der Gegner schlägt vor, C89 (Aufwand 602, mein Wert 389) durch B05 (Aufwand 849, mein Wert 508) zu ersetzen. Mein aktueller Vorschlag (mit C89) hatte einen Aufwand von 41+75+417+129+390+350+118+630+602 = 2752. Mit B05 wäre der Aufwand 2752 {-} 602 + 849 = 2999, was noch unter 3288 liegt. Allerdings ist B05 für mich weniger wertvoll pro Aufwand (508/849 ≈ 0.598) im Vergleich zu C89 (389/602 ≈ 0.646). Der Gesamtwert für mich würde von {[}583+822+868+780+783+461+484+668+389=5838{]} auf {[}5838 {-} 389 + 508 = 5957{]} steigen, aber der Aufwand steigt ebenfalls. Da noch Budget übrig ist (3288 {-} 2999 = 289), könnte ich einen weiteren Gegenstand vorschlagen. Vielleicht A07 (Aufwand 329, mein Wert 138) hinzufügen? Das wäre Aufwand 2999+329=3328, zu hoch. Besser A63 (Aufwand 151, mein Wert 121) hinzufügen? Aufwand 2999+151=3150, Wert 5957+121=6078. Das ist gut für mich. Ich sollte diesen Counter vorschlagen.'\}} \\
\\ 
\texttt{ARGUMENT: \{'Ich verstehe deinen Wunsch, B05 einzubeziehen, aber lass uns auch A63 hinzufügen, um den Gesamtwert zu steigern, während wir komfortabel unter dem Budget bleiben.'\}} \\
\texttt{VORSCHLAG: \{'A36', 'C45', 'B73', 'C13', 'B93', 'B16', 'B33', 'C34', 'B05', 'A63'\}} \\
            }
        }
    }
    & & \\ \\

    \theutterance \stepcounter{utterance}  
    & & \multicolumn{4}{p{0.6\linewidth}}{
        \cellcolor[rgb]{0.9,0.9,0.9}{
            \makecell[{{p{\linewidth}}}]{
                \texttt{\tiny{[GM$\rangle$P2]}}
                \texttt{ARGUMENT: \{'Ich verstehe deinen Wunsch, B05 einzubeziehen, aber lass uns auch A63 hinzufügen, um den Gesamtwert zu steigern, während wir komfortabel unter dem Budget bleiben.'\}} \\
\texttt{VORSCHLAG: \{'A36', 'C45', 'B73', 'C13', 'B93', 'B16', 'B33', 'C34', 'B05', 'A63'\}} \\
            }
        }
    }
    & \\ \\

    \theutterance \stepcounter{utterance}  
    & & & \multicolumn{4}{p{0.6\linewidth}}{
        \cellcolor[rgb]{0.9,0.9,0.9}{
            \makecell[{{p{\linewidth}}}]{
                \texttt{\tiny{[GM$\langle$P2]}}
                \texttt{STRATEGISCHE ÜBERLEGUNGEN: \{'Der Gegner schlägt meinen Vorschlag (9 Gegenstände) plus A63 vor. Aufwand: B05 (849) + A63 (151) = 1000 zusätzlich zu den vorherigen 2150? Nein, mein letzter Vorschlag hatte B05 anstelle von C89, also muss ich neu rechnen: A36(41) + C45(75) + B73(417) + C13(129) + B93(390) + B16(350) + B33(118) + C34(630) + B05(849) = Summe: 41+75=116, +417=533, +129=662, +390=1052, +350=1402, +118=1520, +630=2150, +849=2999. Plus A63(151) = 3150. Das ist unter 3288. Mein Wert: 583 (A36) + 822 (C45) + 868 (B73) + 780 (C13) + 783 (B93) + 461 (B16) + 484 (B33) + 668 (C34) + 508 (B05) = Summe: 583+822=1405, +868=2273, +780=3053, +783=3836, +461=4297, +484=4781, +668=5449, +508=5957. Plus A63(121) = 6078. Das ist ausgezeichnet für mich! Sollte ich annehmen? Der Aufwand ist sicher unter dem Limit, und der Wert ist sehr hoch. Vielleicht kann ich sogar noch C11 hinzufügen? 3150 + 801 = 3951, das ist über dem Limit. Also nein. Dieser Deal ist sehr gut für mich. Ich sollte zustimmen.'\}} \\
\texttt{ARGUMENT: \{'Das ist ein sehr fairer Vorschlag, der einen hohen Gesamtwert für beide Seiten bei komfortabler Einhaltung des Budgets bietet. Ich stimme zu.'\}} \\
\texttt{ZUSTIMMUNG: \{'A36', 'C45', 'B73', 'C13', 'B93', 'B16', 'B33', 'C34', 'B05', 'A63'\}} \\
            }
        }
    }
     \\ \\

    \theutterance \stepcounter{utterance}  
    & & & \multicolumn{2}{p{0.3\linewidth}}{
        \cellcolor[rgb]{0.9,0.9,0.9}{
            \makecell[{{p{\linewidth}}}]{
                \texttt{\tiny{[GM$|$GM]}}
                \texttt{game successful} \\
            }
        }
    }
    & & \\ \\

    \theutterance \stepcounter{utterance}  
    & & & \multicolumn{2}{p{0.3\linewidth}}{
        \cellcolor[rgb]{0.9,0.9,0.9}{
            \makecell[{{p{\linewidth}}}]{
                \texttt{\tiny{[GM$|$GM]}}
                \texttt{end game} \\
            }
        }
    }
    & & \\ \\

\end{supertabular}
}

\end{document}
